\documentclass[twoside]{article}

\usepackage[preprint]{aistats2026}
\usepackage[round,authoryear]{natbib}
\usepackage{amsmath,amssymb,amsthm,mathtools}
\usepackage{booktabs}
\usepackage{tabularx}
\usepackage{microtype}
\usepackage{xcolor}
\usepackage{enumitem}

\renewcommand{\bibname}{References}
\renewcommand{\bibsection}{\subsubsection*{\bibname}}
\bibliographystyle{apalike}

\newtheorem{theorem}{Theorem}
\newtheorem{lemma}[theorem]{Lemma}
\newtheorem{corollary}[theorem]{Corollary}
\newtheorem{proposition}[theorem]{Proposition}
\newtheorem{remark}[theorem]{Remark}
\newtheorem{definition}[theorem]{Definition}

\newcommand{\Hc}{\mathcal H}
\newcommand{\D}{\mathcal D}
\newcommand{\E}{\mathbb E}
\newcommand{\Pp}{\mathbb P}
\newcommand{\one}{\mathbf 1}
\newcommand{\Risk}{R}
\newcommand{\Ex}{\Delta}
\newcommand{\Profile}{\mathsf{Profile}}
\newcommand{\Select}{\mathsf{Select}}
\newcommand{\Tour}{\mathsf{Tournament}}
\newcommand{\poly}{\operatorname{poly}}

\begin{document}
\runningauthor{}

\twocolumn[
\aistatstitle{Exact Adaptation to Unknown Tsybakov Noise}

\aistatsauthor{Pahan Dewasurendra}
\aistatsaddress{Johns Hopkins University}]

\begin{abstract}
The minimax excess-risk rate for a binary class of VC dimension $d$ under Tsybakov noise with exponent $\alpha\in(0,1)$ is
\[
 \left(\frac{d+\log(1/\delta)}{n}\right)^{1/(2-\alpha)}.
\]
A recent proper learner attains this rate for every VC class, but it requires both noise parameters. Whether exact adaptation is possible without either parameter was left open. We answer this question affirmatively. Our learner is proper, receives only the class, sample size, and confidence, and simultaneously attains the minimax rate for every fixed tail constant and exponent. There is no $\log\log n$ or other adaptation loss. Ordinary validation over the required $\Theta(\!\log n)$ noise profiles would incur such a loss. We avoid it through three scale-sensitive devices. First, an off-design analysis of a capped noise-profile learner makes conservative-profile failures decay geometrically with their distance from the true fixed point. Second, a stop-at-first-instability selector charges comparisons to their risk scale instead of the number of profiles. Third, a prior-weighted empirical-Bernstein tournament removes the unknown tail constant while returning one candidate. The key proof is a two-scale concentration argument: deep regional comparisons have large nominal confidence, while shallow comparisons gain variance from the stronger true noise exponent. Their balance makes all profile costs summable. This resolves exact parameter adaptation for arbitrary VC classes while preserving properness and the sharp confidence dependence.
\end{abstract}

\section{INTRODUCTION}

Fast classification rates depend on how much probability lies near the Bayes decision boundary. The Tsybakov condition quantifies this through an exponent $\alpha$ that interpolates between agnostic and nearly noiseless learning \citep{mammen1999smooth,tsybakov2004optimal}. For a binary class $\Hc$ of VC dimension $d$, with the Bayes rule in $\Hc$, the minimax high-probability excess risk is
\begin{equation}
 r_\alpha(n,\delta)
 =\left(\frac{d+\log(1/\delta)}{n}\right)^{1/(2-\alpha)}.
 \label{eq:intro-rate}
\end{equation}
The lower bound has long been known, while empirical risk minimization pays an additional logarithm for a general VC class \citep{massart2006risk,hanneke2015minimax}.

\citet{hanneke2026optimal} recently closed this twenty-year gap. Their MERIT construction isolates regions of increasing noise and intersects regional near-minimizer sets. It is proper and attains \eqref{eq:intro-rate}, but its schedules use the tail constant and $\alpha$. Their first open question asks whether an optimal learner can adapt to both parameters. We resolve that question.

\begin{theorem}[Informal]
There is one proper learner, independent of both Tsybakov parameters, whose excess risk is at most $K(C,\alpha)r_\alpha(n,\delta)$ with probability at least $1-\delta$ for every tail constant $C$, every $\alpha\in(0,1)$, and every VC class.
\end{theorem}

\begin{table*}[t]
\caption{Closest passive-classification guarantees. ``Exact'' means no sample-dependent factor beyond the minimax rate for the stated class. Structured results use boundary or entropy assumptions beyond finite VC dimension.}
\label{tab:comparison}
\centering
\footnotesize
\begin{tabularx}{\textwidth}{@{}lXXll@{}}
\toprule
Work & hypothesis setting & parameters used & rate & output \\
\midrule
\citet{massart2006risk} & arbitrary VC & neither & logarithmic loss & proper ERM \\
\citet{tsybakov2004optimal} & structured entropy, $\rho>0$ & neither margin nor complexity & logarithmic loss & selector \\
\citet{lecue2007simultaneous} & structured entropy, $\rho>0$ & neither margin nor complexity & exact & aggregate \\
\citet{hanneke2026optimal} & arbitrary VC & tail constant and exponent & exact & proper \\
Theorem~\ref{thm:main} & arbitrary VC & neither & exact & proper \\
\bottomrule
\end{tabularx}
\end{table*}

The result is exact in its dependence on $n,d$, and $\delta$. A direct grid over $\alpha$ does not suffice. Constant-factor discretization of the possible fixed-point risks requires $J=\Theta(\!\log n)$ profiles. Assigning confidence $\delta/J$ or validating a finite list replaces $d+\log(1/\delta)$ by $d+\log(1/\delta)+\log\log n$. This is an unbounded loss even when $d$ and $\delta$ are fixed.

Our proof removes this loss by charging every random event to a risk scale. Let $s=(d+\log(1/\delta))/n$. We use the geometric grid $r_j=2^{-j}\sqrt{s}$ and define the exponent $b_j$ by $r_j=s^{1/(2-b_j)}$. If $j_\alpha$ is the last conservative profile with $b_j\leq\alpha$, then $r_{j_\alpha}\asymp r_\alpha$. A profile $\ell$ steps coarser has a true concentration exponent enlarged by
\[
 \left(\frac{r_j}{r_\alpha}\right)^{2-\alpha}
 =2^{(2-\alpha)\ell}.
\]
This geometric gain pays for all conservative profiles without dividing confidence by $J$.

There are two complications. MERIT contains a random number of fresh disagreement tests whose concentration is distribution free. We introduce deterministic, parameter-computable loop caps and allocate confidence over each cap. The extra logarithm is absorbed by the $d\log(1/B)$ term already required at disagreement scale $B$. Final regional comparisons also use internal blocks at many depths. At depth $u$, their exponent has the form
\begin{equation}
 \begin{aligned}
 &cD(u+1)\max\{1,G_\ell 2^{-u}/\poly(u+1)\},\\
 &G_\ell=2^{(2-\alpha)\ell},
 \end{aligned}
 \label{eq:intro-balance}
\end{equation}
where $D=d+\log(1/\delta)$. Above the balance depth, nominal confidence is large. Below it, the true variance is small. This makes the sum over all depths at most $K\exp[-cD(\ell+1)]$.

We then select among the candidates in two steps. For each tail-constant guess, a validation rule scans record-low declared radii and stops at the first instability. False rejections before $j_\alpha$ are summable by geometric radii. After $j_\alpha$, only the first candidate with genuinely excessive risk needs to be detected. Finally, a prior-weighted empirical-Bernstein tournament processes dyadic tail-constant guesses. Every accepted replacement lowers population risk. Failure of the first valid guess to replace the champion yields a Bernstein fixed-point inequality and hence the same oracle rate.

The construction is statistical rather than computational. Like the base arbitrary-VC result, it invokes constrained empirical minimization over general classes. It leaves confidence-independent, expected-risk, and efficient variants open.

\section{SETTING AND RESULT}

Let $(X,Y)\sim\D$ on $\mathcal X\times\{0,1\}$. Write $\eta(x)=\Pp(Y=1\mid X=x)$ and $h^*(x)=\one\{\eta(x)\geq1/2\}$. For $h\in\Hc$, let
\[
 \Risk(h)=\Pp(h(X)\ne Y),\qquad
 \Ex(h)=\Risk(h)-\Risk(h^*).
\]
We assume throughout that $h^*\in\Hc$ and $\operatorname{VC}(\Hc)=d\geq1$.

\begin{definition}[Tsybakov tail class]
\label{def:tn}
For $C\geq1$ and $\alpha\in(0,1)$, write $\D\in\mathrm{TN}(C,\alpha)$ if
\begin{equation}
 \Pp\bigl(|2\eta(X)-1|\leq t\bigr)
 \leq C t^{\alpha/(1-\alpha)},\qquad 0<t\leq1.
 \label{eq:tn}
\end{equation}
\end{definition}

This direct tail parameterization makes nesting transparent. If $b\leq\alpha$, then $\mathrm{TN}(C,\alpha)\subseteq\mathrm{TN}(C,b)$. It also implies the regional Bernstein inequality
\begin{equation}
 \Pp(h\ne h^*,X\in A)
 \leq(C+1)\Ex_A(h)^\alpha,
 \label{eq:regional-bernstein}
\end{equation}
where $\Ex_A(h)=\E[|2\eta(X)-1|\one\{h\ne h^*,X\in A\}]$. Indeed, split at $t=\Ex_A(h)^{1-\alpha}$.

\begin{theorem}[Parameter-free optimal learning]
\label{thm:main}
There is a proper learner $\mathsf A$, taking only $(\Hc,d,n,\delta)$, such that for every $C\geq1$ and $\alpha\in(0,1)$ there is a finite $K(C,\alpha)$ for which, for every $n\geq1$ and $\delta\in(0,1/4)$,
\begin{equation}
 \sup_{\D\in\mathrm{TN}(C,\alpha)}
 \Pp_{S_n\sim\D^n}\!\left(
 \Ex(\mathsf A(S_n,\delta))>K(C,\alpha)r_\alpha(n,\delta)
 \right)
 \leq\delta.
 \label{eq:main}
\end{equation}
The output belongs to $\Hc$, and $\mathsf A$ uses neither $C$ nor $\alpha$.
\end{theorem}

Together with the standard lower bound \citep{massart2006risk,hanneke2015minimax}, Theorem~\ref{thm:main} gives simultaneous minimax optimality. Constants may deteriorate as $\alpha$ approaches one. This is unavoidable for our reduction and is consistent with the qualitative change at the realizable endpoint, where general proper learners cannot attain the optimal rate.

\section{THE PARAMETER-FREE LEARNER}
\label{sec:algorithm}

Split the sample into three constant-size independent blocks $S,V,W$. Profiles are trained on $S$, exponent selection uses $V$, and tail-constant selection uses $W$. Put $D=d+\log(64/\delta)$, $N=\min\{|S|,|V|,|W|\}$, and $s=D/N$. The internal rate $s^{1/(2-\alpha)}$ differs from \eqref{eq:intro-rate} by only a numerical factor. If $s$ is larger than a numerical constant, return an empirical risk minimizer. We henceforth take $s$ small.

\paragraph{Fixed-point profiles.}
Let
\[
 r_j=2^{-j}\sqrt{s},\qquad
 b_j=2-\frac{\log s}{\log r_j},
 \label{eq:grid}
\]
and stop before $r_j<s$. Retain $b_j\in[b_-,1-b_-]$, where $b_-=1/\log\log(1/s)$. An agnostic ERM candidate of declared radius $M_0\sqrt{s}$ covers the omitted lower endpoint. For every fixed $\alpha\in(0,1)$, the grid eventually contains a last $j_\alpha$ with $b_{j_\alpha}\leq\alpha$, and
\begin{equation}
 r_\alpha=s^{1/(2-\alpha)}\leq r_{j_\alpha}<2r_\alpha.
 \label{eq:oracle-grid}
\end{equation}

For tail guesses $C_k=2^k$, $0\leq k\leq\lceil\log_2N\rceil$, run a capped profile learner
\[
 h_{k,j}=\Profile(S,C_k,b_j,\delta).
\]
It is a confidence-strengthened version of MERIT. Appendix~\ref{app:profile} gives the full construction. Its schedules are those of \citet{hanneke2026optimal}, reparameterized by the tail constant in \eqref{eq:tn}, with deterministic disagreement-loop caps and explicit endpoints. Let $M(C_k,b_j)\geq1$ be its declared constant and set
\[
 \rho_{k,j}=M(C_k,b_j)r_j.
\]
If a run exhausts its deterministic sample allocation, it returns the agnostic candidate.

\paragraph{Record thinning and exponent selection.}
For fixed $k$, start with the agnostic candidate and scan profiles from coarse to fine. Keep a profile only when its declared radius is at most half the last kept radius. This yields $(h_0,\rho_0),\ldots,(h_L,\rho_L)$ with $\rho_{q+1}\leq\rho_q/2$.

The selector accepts $h_0$ and scans $q=1,\ldots,L$. It accepts $h_q$ only if
\begin{equation}
 \widehat R_V(h_q)-\widehat R_V(h_p)
 \leq L_0\rho_p\qquad\text{for every }p<q.
 \label{eq:selector}
\end{equation}
At the first rejection it returns the preceding candidate. If no rejection occurs, it returns $h_L$. Denote the result by $g_k$.

\paragraph{Tail-constant tournament.}
Conditionally on $(S,V)$, the $g_k$ are fixed. Give index $k$ weight $\pi_k=6/[\pi^2(k+1)^2]$. For a pair $(k,l)$, define
\begin{align}
 x_{kl}&=\log\frac{32}{\delta\pi_k\pi_l},
 &\widehat p_{kl}&=\widehat\Pp_W(g_k\ne g_l),\nonumber\\
 U_{kl}&=2\widehat p_{kl}+3x_{kl}/|W|,
 &q_{kl}&=x_{kl}/|W|,\nonumber\\
 \operatorname{rad}_{kl}
 &=\sqrt{2U_{kl}q_{kl}}+\frac{2q_{kl}}{3}.
 \label{eq:pair-radius}
\end{align}
Start with champion $c=0$. Challenger $k$ replaces $c$ if
\begin{equation}
 \widehat R_W(g_k)-\widehat R_W(g_c)
 +\operatorname{rad}_{kc}<0.
 \label{eq:tournament}
\end{equation}
Return the final champion. This completes $\mathsf A$.

\section{WHY ADAPTATION COSTS NO LOGARITHM}
\label{sec:proof}

We state the three ingredients. Proofs, constants, and the profile event table appear in the appendix.

\begin{lemma}[Off-design profile]
\label{lem:profile}
Fix a true $\D\in\mathrm{TN}(C,\alpha)$ and a valid guess $C_k\geq C$. With probability at least $1-\delta/16$, every conservative profile $b_j\leq\alpha$ obeys
\[
 \Ex(h_{k,j})\leq\rho_{k,j}.
\]
Moreover, the failure contribution of the profile $\ell=j_\alpha-j$ scales above the oracle is at most
\[
 K(C,\alpha)\exp[-c(C,\alpha)D(\ell+1)].
\]
\end{lemma}

The core calculations are visible from the schedules. At internal depth $u$, write $p_b=(1-b)/(2-b)$. Up to constants and powers of $u+1$,
\begin{align}
 n_t&=N2^{-u},
 &a_t&=r_b^{-(1-b)}2^{-u(1+p_b)},\nonumber\\
 \gamma_t&=r_b2^{u/(2-b)}(u+1)^{1/(2-b)},
 &B_t&=a_t\gamma_t,
 \label{eq:schedules}\\
 e_t&=r_b2^{-p_bu}(u+1).\nonumber
\end{align}
The pruning exponent under the true $\alpha$ is exactly, up to a fixed factor,
\begin{equation}
 D\left(\frac{r_b}{r_\alpha}\right)^{2-\alpha}
 2^{-u(\alpha-b)/(2-b)}
 (u+1)^{(2-\alpha)/(2-b)}.
 \label{eq:pruning-exp}
\end{equation}
One grid step changes $b$ by $O(1/\log(1/s))$, while $u=O(\log\log(1/s))$. Thus the second factor cannot cancel the geometric fixed-point gain.

For final regional comparisons, the disagreement variance is bounded both by the design scale $B_t$ and by $K\gamma_t^\alpha$ from \eqref{eq:regional-bernstein}. Their ratio produces \eqref{eq:intro-balance}. Splitting the depths at the balance point gives a geometric sum rather than a factor $\log\log n$. The terminal region has exponent $Nr_b^{2-\alpha}$ directly.

The disagreement tests do not enjoy a stronger-noise exponent, so we make them simultaneous. The number of successful updates has the computable cap
\[
 S_t^{\max}=\left\lceil
 K(C_k,b)a_t^{-1/(1-b)}\gamma_t^{-1}
 \right\rceil.
\]
Charging each of these draws and the final unsuccessful draw adds
\[
 O_b(\log S_t^{\max}+\log J+\log T)
 =O_b(\log\log(1/s))
\]
to its confidence logarithm. Every retained profile has $\log(1/B_t)\gtrsim\log(1/s)/\log\log(1/s)$, so the existing relative-VC sample size absorbs this addition. Lemma~\ref{lem:profile} follows by the localized VC inequality and a sum over \eqref{eq:pruning-exp} and \eqref{eq:intro-balance}.

Record thinning preserves an oracle profile. Indeed, $M(C_k,b)$ is bounded on $b\in[0,\alpha]$ for fixed $(C_k,\alpha)$. The last record at or before $j_\alpha$ has radius at most $2\rho_{k,j_\alpha}\leq K(C,\alpha)r_\alpha$.

\begin{lemma}[Ordered scale selection]
\label{lem:selector}
Suppose $\rho_{q+1}\leq\rho_q/2$ and every candidate through $q^*$ satisfies $\Ex(h_q)\leq\rho_q$, with $\rho_{q^*}\leq Kr_\alpha$. Under \eqref{eq:regional-bernstein}, the rule \eqref{eq:selector} returns $g$ satisfying $\Ex(g)\leq K'r_\alpha$ except with probability at most
\[
 K'\exp[-c|V|r_\alpha^{2-\alpha}].
\]
\end{lemma}

For $p<q\leq q^*$, a false rejection needs an upward deviation of order $\rho_p$. Its probability is at most $\exp[-c|V|\rho_p^{2-\alpha}]$. At most $q^*-p$ comparisons use scale $\rho_p$, and
\begin{equation}
 \sum_{l\geq0}(l+1)
 \exp[-c|V|r_\alpha^{2-\alpha}2^{(2-\alpha)l}]
 \leq K\exp[-c'|V|r_\alpha^{2-\alpha}].
 \label{eq:selector-sum}
\end{equation}
After $q^*$, all candidates are harmless until the first one whose excess exceeds a large multiple of $r_\alpha$. Acceptance of that candidate requires a downward deviation in its comparison with $h_{q^*}$. Stopping at the first rejection avoids a union over all later profiles.

\begin{lemma}[Prior-weighted tournament]
\label{lem:tournament}
There is an event of probability at least $1-\delta/8$ with the following property. For each candidate pair, put
\[
 E_{kl}=(\widehat R_W-R)(g_k)
 -(\widehat R_W-R)(g_l).
\]
On this event, simultaneously for all pairs,
\begin{align}
 |E_{kl}|&\leq\operatorname{rad}_{kl},\nonumber\\
 \operatorname{rad}_{kl}
 &\leq3\sqrt{\Pp(g_k\ne g_l)q_{kl}}
 +5q_{kl}.
 \label{eq:empbern}
\end{align}
Every replacement in \eqref{eq:tournament} lowers population risk. If candidate $g_k$ does not replace the current champion $g_c$, then
\[
 \Ex(g_c)-\Ex(g_k)\leq2\operatorname{rad}_{ck}.
\]
\end{lemma}

The proof is ordinary Bernstein plus a binomial upper bound for the disagreement probability. The summable prior permits a union over the growing list.

\paragraph{Proof of Theorem~\ref{thm:main}.}
Fix $C$ and $\alpha$, and let $k^*=\lceil\log_2C\rceil$. For all sufficiently large $n$, $k^*$ is on the outer grid. Lemmas~\ref{lem:profile} and \ref{lem:selector} give
\[
 \Ex(g_{k^*})\leq K(C,\alpha)r_\alpha
\]
with the required probability. Immediately before challenge $k^*$, the champion index is below $k^*$. If $g_{k^*}$ replaces it, the champion has oracle risk. Otherwise Lemma~\ref{lem:tournament} and \eqref{eq:regional-bernstein} give, with $q=O_C(\log(1/\delta)/n)$,
\begin{equation}
 \Ex(g_c)-\Ex(g_{k^*})
 \leq K\sqrt{(\Ex(g_c)^\alpha+\Ex(g_{k^*})^\alpha)q}+Kq.
 \label{eq:fixedpoint}
\end{equation}
If $\Ex(g_c)\leq2\Ex(g_{k^*})$, the claim follows. Otherwise solving \eqref{eq:fixedpoint} yields
\[
 \Ex(g_c)\leq
 K\Ex(g_{k^*})+Kq^{1/(2-\alpha)}+Kq.
\]
All later replacements lower risk. The profile and selector constants are fixed so their failure probabilities are at most $\delta/16$ each. Together with the $\delta/8$ tournament event, a union bound proves \eqref{eq:main} for large $n$. Enlarging $K(C,\alpha)$ covers the remaining finite range. Properness follows because every stage returns a member of $\Hc$. \hfill$\square$

\section{RELATED WORK AND LIMITS}

The original margin and aggregation theory established fast rates and robustness to an unknown margin for structured boundary classes \citep{mammen1999smooth,tsybakov2004optimal}. The adaptive theorem of \citet{tsybakov2004optimal} uses nested finite sieves with bracketing exponent in a fixed subinterval of $(0,1)$ and retains powers of $\log n$. The square-root penalty of \citet{tsybakov2005square} also adapts only up to logarithms.

\citet{lecue2007simultaneous} obtains exact simultaneous adaptation when the bracketing exponent $\rho$ stays in a compact subset of $(0,1)$. Its aggregation remainder is lower order because $\rho>0$. An arbitrary VC class is the zero-entropy-exponent boundary, where that remainder has the same polynomial order as the target and its logarithm cannot be absorbed. Our scale-charged selector addresses precisely this boundary. Local empirical-process tools behind the profile analysis originate in ratio and local-complexity inequalities \citep{gine2006concentration,koltchinskii2006local,bartlett2006empirical}.

Strong margin adaptation can be impossible for arbitrary nonnested model collections with model-specific local margins \citep{arlot2011margin}. That negative result does not conflict with Theorem~\ref{thm:main}. Here one global tail condition controls every candidate, and the profiles form a specific nested difficulty scale. Online methods can adapt to unknown Bernstein exponents, but standard learning-rate mixtures carry second-order terms such as $\log\log n$ \citep{koolen2016combining}. Our ordered rule is tailored to proper batch classification and avoids that term.

This work does not make MERIT efficient. It also remains confidence dependent. A confidence-free learner would imply an optimal expected-risk result, which is another open question from \citet{hanneke2026optimal}. Adaptation at the endpoint $\alpha=1$ is excluded because arbitrary-class proper learning changes character there.

\section{CONCLUSION}

Exact arbitrary-VC learning under Tsybakov noise does not require the noise parameters. The apparent model-selection cost disappears when training failures, validation comparisons, and constant guesses are each charged at their natural variance scale. The resulting learner is proper and preserves the minimax dependence on sample size, dimension, and confidence.

\bibliography{references}

\appendix
\section{PROOFS FOR THE PROFILE LEARNER}
\label{app:profile}

This appendix gives the capped profile construction and proves Lemma~\ref{lem:profile}. Write $P$ for probability under $\D$ and $\ell_h(x,y)=\one\{h(x)\ne y\}$. Numerical constants are denoted by $c,C$, while $K(C_0,b)$ may depend on a declared tail constant and exponent. Their values can change from line to line. All logarithms below are at least one. Floors and ceilings only change numerical constants.

\subsection{Concentration and tail consequences}

For a measurable region $A$, define
\[
 \Ex_A(h)=\E\!\left[|2\eta(X)-1|
 \one\{h(X)\ne h^*(X),X\in A\}\right].
\]

\begin{lemma}[Regional Bernstein condition]
\label{lem:regional-appendix}
If $\D\in\mathrm{TN}(C,\alpha)$, then for every $A$ and $h$,
\[
 \Pp(h\ne h^*,X\in A)\leq(C+1)\Ex_A(h)^\alpha.
\]
If $b\leq\alpha$, then $\D\in\mathrm{TN}(C,b)$.
\end{lemma}

\begin{proof}
The nesting follows because $t^{\alpha/(1-\alpha)}\leq t^{b/(1-b)}$ on $[0,1]$. Put $z=\Ex_A(h)^{1-\alpha}$. Splitting the disagreement event according to whether $|2\eta-1|\leq z$ gives
\begin{align*}
 \Pp(h\ne h^*,X\in A)
 &\leq Cz^{\alpha/(1-\alpha)}+\frac{\Ex_A(h)}{z}\\
 &=(C+1)\Ex_A(h)^\alpha.
\end{align*}
The cases $\Ex_A(h)=0$ and $\Ex_A(h)>1$ are immediate.
\end{proof}

We use the following standard localized VC inequality. It follows by peeling a relative empirical-process inequality for the loss-difference class, whose VC dimension is at most a numerical multiple of $d$ \citep{gine2006concentration,koltchinskii2006local}.

\begin{lemma}[Localized VC comparison]
\label{lem:local-vc}
Let $A$ be independent of an i.i.d. sample of size $m$. With probability at least $1-2e^{-x}$, simultaneously for $f,g\in\Hc$,
\begin{align}
 &\left|(P_m-P)
 \bigl[(\one\{f(X)\ne Y\}-\one\{g(X)\ne Y\})\one\{X\in A\}\bigr]\right|
 \nonumber\\
 &\quad\leq C\left[
 \sqrt{\frac{p_{fg,A}}{m}
 \left(d\log Q_{fg,A}+x\right)}
 +\frac{d\log Q_{fg,A}+x}{m}
 \right],
 \label{eq:local-vc}
\end{align}
where $p_{fg,A}=P(f\ne g,X\in A)$ and
\[
 Q_{fg,A}=\frac{eP(A)}{p_{fg,A}}
 \wedge\frac{emP(A)}d.
\]
Logarithms are interpreted as one when their argument is smaller than $e$.
\end{lemma}

We also need a relative disagreement estimate. A standard relative VC bound applied to $\{\one\{f\ne g\}\one_A:f,g\in\Hc\}$ gives the next statement.

\begin{lemma}[Multiplicative disagreement test]
\label{lem:relative-vc}
There is a numerical $C$ such that, if
\[
 m\geq\frac{C}{B}\left(d\log\frac1B+x\right),
\]
then, with probability at least $1-2e^{-x}$, simultaneously for $f,g\in\Hc$ and every fixed $A$,
\begin{align*}
 P_m(f\ne g,X\in A)\geq B&\ \Longrightarrow\
 P(f\ne g,X\in A)\geq B/2,\\
 P_m(f\ne g,X\in A)\leq B&\ \Longrightarrow\
 P(f\ne g,X\in A)\leq2B.
\end{align*}
The result remains valid conditionally when $A$ is determined by earlier independent blocks.
\end{lemma}

The next deterministic lemma controls the mass isolated by a disagreement loop. It is the population step behind MERIT.

\begin{lemma}[Noisy-region mass]
\label{lem:region-mass}
Suppose $V\subseteq\Hc$ and $\Ex_A(h)\leq\gamma$ for every $h\in V$. Starting from $U_0=\varnothing$, let
\[
 U_q=U_{q-1}\cup\{f_q\ne g_q\},
 \qquad f_q,g_q\in V,
\]
and assume $P(f_q\ne g_q,X\in A\setminus U_{q-1})\geq B$. If $\D\in\mathrm{TN}(C,\alpha)$, then
\[
 P(U_q\cap A)\leq K(C,\alpha)(\gamma/B)^{\alpha/(1-\alpha)}.
\]
Consequently $q\leq K(C,\alpha)B^{-1}(\gamma/B)^{\alpha/(1-\alpha)}$.
\end{lemma}

\begin{proof}
The pieces $Q_i=\{f_i\ne g_i\}\cap(A\setminus U_{i-1})$ are disjoint. On the disagreement of two binary classifiers, one of them disagrees with $h^*$, so
\[
 \E[|2\eta-1|\one_{Q_i}]
 \leq\Ex_A(f_i)+\Ex_A(g_i)\leq2\gamma.
\]
Since $P(Q_i)\geq B$, the average margin on each piece is at most $2\gamma/B$. The same holds on their union $U_q\cap A$. Put $p=P(U_q\cap A)$ and $z=(p/(2C))^{(1-\alpha)/\alpha}$, truncated at one. The tail condition gives
\[
 \E[|2\eta-1|\mid U_q\cap A]
 \geq z\left(1-\frac{Cz^{\alpha/(1-\alpha)}}p\right)
 =z/2.
\]
Comparison with $2\gamma/B$ gives the mass bound. The pieces have mass at least $B$, which proves the count bound.
\end{proof}

\subsection{Capped profile construction}

Fix a declared $(C_0,b)$ with $C_0\geq1$ and $b\in(0,1)$. Let $m$ be a fixed constant fraction of $|S|$, set
\[
 D=d+\log(64/\delta),\qquad s=D/m,
 \qquad r_b=s^{1/(2-b)},
\]
and assume $s$ is smaller than a numerical constant. Put
\begin{align*}
 T&=\left\lceil\log_2\log_2(1/s)\right\rceil+1,\\
 n_t&=m2^{t-T},\qquad 0\leq t\leq T.
\end{align*}
The endpoint $n_0$ is defined even though no pruning block uses it. Let $p_b=(1-b)/(2-b)$ and define, for $0\leq t\leq T$,
\begin{equation}
 a_t=2^{t-T}\left(\frac{n_t}{D}\right)^{p_b}.
 \label{eq:appendix-a}
\end{equation}
For $1\leq t\leq T$, put $\zeta_t=(\delta/128)^{T+1-t}$ and
\begin{align}
 \Lambda_t={}&d\log\!\left[
 \left(\frac{n_t}{D}\right)^{b/(2-b)}
 a_{t-1}^{-b/(1-b)}\right]
 +\log(2/\zeta_t),
 \label{eq:appendix-Lambda}\\
 \gamma_t&=A_1(C_0,b)
 \left(\frac{\Lambda_t}{n_t}\right)^{1/(2-b)},
 \label{eq:appendix-gamma}\\
 B_t&=a_t\gamma_t.
 \label{eq:appendix-B}
\end{align}
The log is at least one. For $0\leq t<T$, define
\begin{equation}
 e_t=A_2(C_0,b)2^{-p_b(T-t)}(T-t+1)r_b,
 \label{eq:appendix-e}
\end{equation}
and set $e_T=A_2(C_0,b)r_b$. The constants $A_1,A_2$ are chosen large enough for the comparison events below.

The profile reuses fixed portions of $S$ as follows. One portion supplies disjoint pruning blocks $S_{1,t}$ of size $n_t$. One portion, of size $m$, is the common regional block $S_3$. A third portion supplies fresh disagreement blocks. The same portions may be reused across different profiles, since the later selectors use independent $V$ and all profile probabilities are controlled jointly.

At stage $t=1,\ldots,T$, let $\widehat h^*_{t-1}$ minimize empirical risk on $S_{1,t}\cap\Delta_{t-1}$ and define
\begin{align}
 \Hc_{t-1}=\bigl\{h\in\Hc:
 &\widehat R_{S_{1,t}}(h,\Delta_{t-1})\nonumber\\
 &\leq\widehat R_{S_{1,t}}
 (\widehat h^*_{t-1},\Delta_{t-1})+\gamma_t\bigr\}.
 \label{eq:pruned-class}
\end{align}
Initialize $\Delta_t=\varnothing$. Let
\begin{equation}
 L_t=\left\lceil
 K_0(C_0,b)a_t^{-1/(1-b)}\gamma_t^{-1}
 \right\rceil+1.
 \label{eq:loop-cap}
\end{equation}
For each of at most $L_t$ iterations, draw a fresh block of size
\begin{equation}
 m_t=\left\lceil\frac{A_3}{B_t}
 \left[d\log\frac1{B_t}
 +\log\frac{32JT L_t}{\zeta_t}\right]\right\rceil.
 \label{eq:disagreement-size}
\end{equation}
If there are $f,g\in\Hc_{t-1}$ with empirical disagreement at least $B_t$ on $\Delta_{t-1}\setminus\Delta_t$, update $\Delta_t\leftarrow\Delta_t\cup\{f\ne g\}$. Otherwise terminate the loop. If it has not terminated by iteration $L_t$, or if the cumulative fresh-sample allocation exceeds its portion of $S$, abort the profile.

For $A_t=\Delta_t\setminus\Delta_{t+1}$, with $\Delta_0=\mathcal X$, let $\widehat h_t$ minimize empirical regional risk over $\Hc_t$ on $S_3\cap A_t$ and set
\begin{align}
 \widehat\Hc_t=\{h\in\Hc_t:
 &\widehat R_{S_3}(h,A_t)\nonumber\\
 &\leq\widehat R_{S_3}(\widehat h_t,A_t)+e_t\},
 \qquad 0\leq t<T.
 \label{eq:regional-set}
\end{align}
On the terminal region, minimize over $\Hc$ and define $\widehat\Hc_T$ with tolerance $A_2r_b$. If the profile has not aborted and $\cap_{t=0}^T\widehat\Hc_t$ is nonempty, output any member. Otherwise output the agnostic ERM candidate. This is $\Profile(S,C_0,b,\delta)$.

The construction is a capped and confidence-strengthened form of the proper learner in \citet{hanneke2026optimal}. Equations \eqref{eq:appendix-a}--\eqref{eq:appendix-e} retain its essential schedules. We explicitly define $a_0$, charge the final unsuccessful disagreement draw, and use the direct tail constant of Definition~\ref{def:tn}.

\subsection{Schedule identities and sample budget}

Let $u=T-t$. Direct substitution gives the following bounds, uniformly for $0\leq u\leq T$, with constants depending only on $(C_0,b)$:
\begin{align}
 a_t&=r_b^{-(1-b)}2^{-u(1+p_b)},
 \label{eq:a-identity}\\
 \gamma_t&\asymp r_b2^{u/(2-b)}(u+1)^{1/(2-b)},
 \label{eq:gamma-identity}\\
 B_t&\asymp r_b^b2^{-2p_bu}(u+1)^{1/(2-b)},
 \label{eq:B-identity}\\
 e_t&\asymp r_b2^{-p_bu}(u+1).
 \label{eq:e-identity}
\end{align}
For \eqref{eq:gamma-identity}, the logarithm in \eqref{eq:appendix-gamma} is $O_b(u+1)$ and $\log(1/\zeta_t)\asymp(u+1)\log(1/\delta)$. Increasing $A_1$ supplies matching lower bounds.

The loop cap satisfies
\begin{equation}
 \log L_t\leq K(C_0,b)+\frac{u\log2}{1-b}+\log(u+1).
 \label{eq:loop-log}
\end{equation}
Indeed, use \eqref{eq:a-identity} and \eqref{eq:gamma-identity} in \eqref{eq:loop-cap}. The powers of $r_b$ cancel.

Let $L=\log(1/s)$. The exponent grid has $J=O(L)$ profiles and $T=O(\log L)$. Every retained $b\geq1/\log L$ obeys, by \eqref{eq:B-identity},
\begin{equation}
 \log(1/B_t)\geq c\frac{L}{\log L}
 \label{eq:B-log-lower}
\end{equation}
for all sufficiently large $L$. For every fixed true $\alpha<1$, \eqref{eq:loop-log} is $O_\alpha(\log L)$ uniformly over $b\leq\alpha$. Thus the extra confidence term in \eqref{eq:disagreement-size} is absorbed by a constant multiple of $d\log(1/B_t)+\log(1/\zeta_t)$.

The known-profile sample calculation then applies with a constant enlargement. For completeness, substitution of \eqref{eq:a-identity}--\eqref{eq:B-identity} into $L_tm_t$ gives a fixed power of $2^u$ times
\[
 m^{b/(2-b)}D^{(2-2b)/(2-b)}\poly(\log(m/D)).
\]
Summing $u\leq T$ contributes another fixed power of $\log(m/D)$. Since $(2-2b)/(2-b)>0$ for every fixed $b<1$, the ratio of this quantity to $m$ tends to zero. The convergence is uniform on $b\in[b_-,\alpha]$ for fixed $\alpha<1$. Therefore every conservative profile stays inside its deterministic sample allocation for all sufficiently large $m$. The cap makes the algorithm well defined for every sample.

\subsection{Off-design event probabilities}

Fix now a true $\D\in\mathrm{TN}(C,\alpha)$, a guess $C_0\geq C$, and a conservative $b\leq\alpha$. Let $r_\alpha=s^{1/(2-\alpha)}$ and write $r_b/r_\alpha=2^\ell$, allowing a displacement of at most one due to the grid endpoint.

All disagreement implications hold simultaneously over profiles, stages, and fresh draws except with probability at most $\delta/128$. This follows from Lemma~\ref{lem:relative-vc}, \eqref{eq:disagreement-size}, and the union over the known caps. On this event, Lemma~\ref{lem:region-mass} and induction give
\begin{align}
 P(\Delta_t)&\leq K(C,\alpha)a_t^{-\alpha/(1-\alpha)},
 \label{eq:true-region-mass}\\
 P(f\ne g,X\in\Delta_{t-1}\setminus\Delta_t)
 &\leq2B_t,\qquad f,g\in\Hc_{t-1}.
 \label{eq:design-disagreement}
\end{align}
The true region bound is at least as strong as the design bound for $a_t\geq1$.

The remaining events have the following scales.
\begin{table}[h]
\caption{Off-design comparison events. Constants and powers of $u+1$ are shown in the subsequent displays.}
\label{tab:event-scales}
\centering
\small
\begin{tabular}{llll}
\toprule
event & sample & threshold & variance proxy \\
\midrule
pruning & $n_t$ & $\gamma_t$ & $K\gamma_t^\alpha$ \\
regional & $m$ & $e_t$ & $K\min\{B_{t+1},\gamma_{t+1}^\alpha\}$ \\
terminal & $m$ & $r_b$ & $K P(\Delta_T)$ \\
\bottomrule
\end{tabular}
\end{table}

We make the event induction explicit. Conditional on the preceding regions, apply Lemma~\ref{lem:local-vc} to the pruning block with $g=h^*$. The empirical near-minimizer definition and the two comparison inequalities, one for $h$ and one for the empirical minimizer, imply on this event
\begin{equation}
 h^*\in\Hc_{t-1},\qquad
 \sup_{h\in\Hc_{t-1}}\Ex_{\Delta_{t-1}}(h)
 \leq K(C,\alpha)\gamma_t.
 \label{eq:pruning-transfer}
\end{equation}
To verify the second statement, if the excess $z$ were larger than $K\gamma_t$, then \eqref{eq:regional-bernstein} and Lemma~\ref{lem:local-vc} would make the empirical excess at least
\[
 z-C\sqrt{z^\alpha\Lambda_t/n_t}-C\Lambda_t/n_t,
\]
where $\Lambda_t$ is the localized entropy plus confidence. The schedule gives $\Lambda_t/n_t\leq K\gamma_t^{2-b}$. Since $b\leq\alpha$ and $\gamma_t\leq1$, this is at most $K\gamma_t^{2-\alpha}$. For sufficiently large $K$, the last display exceeds $\gamma_t$, contradicting membership in \eqref{eq:pruned-class}. The same calculation with the empirical minimizer proves that $h^*$ is retained.

Given \eqref{eq:pruning-transfer}, the simultaneous disagreement event and Lemma~\ref{lem:region-mass} prove \eqref{eq:true-region-mass}. Apply the same lemma with the declared pair $(C_0,b)$ to control the loop length. Every successful empirical update adds at least $B_t/2$ population mass, so the number of updates is strictly below \eqref{eq:loop-cap} when $K_0(C_0,b)$ is large enough. Thus the final unsuccessful test occurs before the cap, and its second implication proves \eqref{eq:design-disagreement}. This closes the region-construction induction.

The regions depend only on pruning and disagreement blocks, so they are independent of $S_3$. Conditional on all regions, apply Lemma~\ref{lem:local-vc} to the loss differences on $A_t$. On the regional comparison event, the ERM argument just used gives
\begin{equation}
 h^*\in\widehat\Hc_t,\qquad
 \sup_{h\in\widehat\Hc_t}\Ex_{A_t}(h)\leq2e_t,
 \quad 0\leq t<T.
 \label{eq:regional-transfer}
\end{equation}
The terminal event similarly gives $h^*\in\widehat\Hc_T$ and terminal excess at most $K r_b$. Hence the intersection is nonempty, and any member has excess bounded by the sum of the regional tolerances and the terminal tolerance.

For pruning, Lemma~\ref{lem:local-vc} and \eqref{eq:true-region-mass} show that failure has exponent at least a fixed multiple of $n_t\gamma_t^{2-\alpha}$. The localized entropy term is no larger than $K_\alpha d(u+1)$. To see this, its potentially growing part is bounded by
\[
 \left(\frac{n_t}{D}\right)^{\alpha/(2-\alpha)}
 a_{t-1}^{-\alpha/(1-\alpha)}.
\]
The exponent of $n_t/D$ is nonpositive when $b\leq\alpha$, since $(1-b)/(2-b)\geq(1-\alpha)/(2-\alpha)$. Only powers of $2^u$ remain.

Using \eqref{eq:gamma-identity} and the fixed-point identities $mr_b^{2-b}=mr_\alpha^{2-\alpha}=D$ gives
\begin{equation}
 n_t\gamma_t^{2-\alpha}
 \geq cD
 \left(\frac{r_b}{r_\alpha}\right)^{2-\alpha}
 2^{-u(\alpha-b)/(2-b)}
 (u+1)^{(2-\alpha)/(2-b)}.
 \label{eq:exact-pruning}
\end{equation}

We record the grid calculation used to lower bound this and the regional events.

\begin{lemma}[Grid-distance bounds]
\label{lem:grid-distance}
There are numerical $c,C>0$ such that, for every fixed $\alpha<1$ and all sufficiently small $s$, if $b\leq\alpha$ is $\ell$ risk-grid steps coarser than the last conservative point, put $G=(r_b/r_\alpha)^{2-\alpha}$. Then, uniformly for $0\leq u\leq T$,
\begin{equation}
 G2^{-u(\alpha-b)/(2-b)}\geq c(\ell+1).
 \label{eq:grid-prune}
\end{equation}
Define
\begin{align}
 \Phi_u=(u+1)^{1+p_b}\max\bigl\{1,
 &G2^{-u(1+(\alpha-b)/(2-b))}\nonumber\\
 &\mathrel{\phantom{G}}\cdot
 (u+1)^{(1-\alpha)/(2-b)}\bigr\}.
 \label{eq:grid-phi}
\end{align}
Then
\begin{equation}
 \sum_{u=0}^T e^{-cD\Phi_u}
 \leq K_\alpha e^{-c_\alpha D(\ell+1)}.
 \label{eq:grid-regional}
\end{equation}
\end{lemma}

\begin{proof}
Write $L=\log_2(1/s)$. On the continuous grid, $-\log_2r_b=L/(2-b)$. The derivative of $b=2-L/(-\log_2r_b)$ with respect to $-\log_2r_b$ is at most $4/L$. Hence
\[
 0\leq\alpha-b\leq\frac{4(\ell+1)}L,
 \qquad
 \frac{r_b}{r_\alpha}\geq2^{\ell-1}.
\]
Also $T\leq C\log L$. The correction in \eqref{eq:grid-prune} is therefore at least
\[
 2^{-C(\ell+1)\log L/L},
\]
which is dominated by $2^{(2-\alpha)(\ell-1)}$. This proves \eqref{eq:grid-prune}, after changing constants for bounded $L$.

For \eqref{eq:grid-regional}, the same correction reduces the fixed-point gain by at most $2^{o(\ell+1)}$. It is enough to bound
\[
 \sum_{u=0}^T
 \exp\{-cD(u+1)\max(1,2^{c_0\ell-u})\}.
\]
Split at $u=c_0\ell/2$. Below the split, the second term is exponentially large in the distance to the split. Above it, the nominal factor $u+1$ is at least a constant multiple of $\ell+1$. The two resulting geometric sums prove the claim.
\end{proof}

By \eqref{eq:exact-pruning} and Lemma~\ref{lem:grid-distance}, the pruning failures for this profile sum to at most $K_\alpha e^{-c_\alpha D(\ell+1)}$.

For a regional comparison, \eqref{eq:design-disagreement} gives variance at most $KB_{t+1}$. The pruning event and Lemma~\ref{lem:regional-appendix} also give variance at most $K\gamma_{t+1}^\alpha$. From \eqref{eq:B-identity}--\eqref{eq:e-identity}, the nominal exponent $me_t^2/B_{t+1}$ is at least $cD(u+1)^{1+p_b}$. The ratio between the design and true variance proxies is, up to fixed constants,
\begin{equation}
 \left(\frac{r_b}{r_\alpha}\right)^{2-\alpha}
 2^{-u(1+(\alpha-b)/(2-b))}
 (u+1)^{(1-\alpha)/(2-b)}.
 \label{eq:variance-ratio}
\end{equation}
The entropy terms remain of order $d(u+1)$. For the design variance, this is the same cancellation used in the known-profile schedule. For the true variance, use \eqref{eq:true-region-mass} to bound the relevant ratio by
\begin{equation}
 \frac{a_t^{-\alpha/(1-\alpha)}}{\gamma_{t+1}^\alpha}
 \leq K_\alpha
 r_b^{\alpha(\alpha-b)/(1-\alpha)}2^{K_\alpha u}
 \leq K_\alpha2^{K_\alpha u}.
 \label{eq:true-entropy-ratio}
\end{equation}
Thus its logarithm is at most $K_\alpha(u+1)$. Lemma~\ref{lem:local-vc}, \eqref{eq:variance-ratio}, and \eqref{eq:grid-regional} now bound the sum of all regional failures by $K_\alpha e^{-c_\alpha D(\ell+1)}$.

Finally, \eqref{eq:true-region-mass} at $T$ gives $P(\Delta_T)\leq K r_b^\alpha$. The envelope form of Lemma~\ref{lem:local-vc} implies, uniformly over terminal comparisons,
\begin{align*}
 &|(P_m-P)(\ell_f-\ell_g)\one_{\Delta_T})|\\
 &\quad\leq K\left[\sqrt{\frac{P(\Delta_T)(d+x)}m}
 +\frac{d+x}m\right].
\end{align*}
Here we used $p\log(eP(\Delta_T)/p)\leq P(\Delta_T)$ and the truncation in $Q_{fg,A}$ to absorb the additive entropy term. Taking $x=cD(r_b/r_\alpha)^{2-\alpha}$ makes the display at most $K r_b$ and gives failure probability
\[
 2\exp\{-cD(r_b/r_\alpha)^{2-\alpha}\},
\]
which is stronger than required.

On the complement of these failures, $h^*$ belongs to every empirical regional set. Every member of their intersection has regional excess at most $2e_t$ on stage $t$ and at most $K r_b$ on the terminal region. Since
\[
 \sum_{u\geq1}2^{-p_bu}(u+1)<\infty,
\]
the profile output has excess at most $M(C_0,b)r_b$ for an explicit finite $M$. The constants are continuous and bounded for $b$ in every compact subset of $[0,1)$. Summing the geometric profile failures over $\ell$ and adding the simultaneous disagreement event proves Lemma~\ref{lem:profile}.

\section{ORDERED EXPONENT SELECTION}
\label{app:selector}

We prove Lemma~\ref{lem:selector} in a form that allows dependence among the trained candidates. Condition on the training block, so all candidates and radii are fixed before $V$.

\begin{proof}[Proof of Lemma~\ref{lem:selector}]
Let $q^*$ be the oracle record. For $p<q\leq q^*$, the population risk difference obeys
\[
 \Ex(h_q)-\Ex(h_p)\leq\Ex(h_q)
 \leq\rho_q\leq\rho_p/2.
\]
Choose $L_0$ larger than a numerical constant. A false rejection requires an upward deviation of at least $cL_0\rho_p$. The loss difference has variance
\begin{align*}
 P(h_q\ne h_p)
 &\leq P(h_q\ne h^*)+P(h_p\ne h^*) \\
 &\leq K(C,\alpha)(\rho_q^\alpha+\rho_p^\alpha)
 \leq K\rho_p^\alpha.
\end{align*}
Bernstein's inequality makes its probability at most $2\exp[-c|V|\rho_p^{2-\alpha}]$. At most $q^*-p$ comparisons use reference scale $\rho_p$. Since $\rho_p\geq2^{q^*-p}\rho_{q^*}$, their complete sum is
\[
 2\sum_{l\geq0}(l+1)
 \exp[-c|V|\rho_{q^*}^{2-\alpha}2^{(2-\alpha)l}],
\]
which is bounded by $K\exp[-c'|V|\rho_{q^*}^{2-\alpha}]$ whenever the base exponent is at least a numerical constant.

Now consider candidates after $q^*$. If all have excess at most $K_1\rho_{q^*}$, every possible output is good. Otherwise let $q_{\mathrm{bad}}$ be the first candidate with excess larger than $K_1\rho_{q^*}$. Any earlier post-oracle candidate is good. Acceptance of $q_{\mathrm{bad}}$ requires its comparison with $h_{q^*}$ to deviate downward by a fixed multiple of $\Ex(h_{q_{\mathrm{bad}}})$. Its variance is at most
\[
 K(\Ex(h_{q_{\mathrm{bad}}})^\alpha+\rho_{q^*}^\alpha),
\]
so Bernstein bounds the acceptance probability by $2\exp[-c|V|\rho_{q^*}^{2-\alpha}]$. A rejection before this index returns a good candidate. Rejection at this index also returns a good candidate. The rule stops there, so no later union is needed.

Finally $|V|\rho_{q^*}^{2-\alpha}\geq cD$ because $\rho_{q^*}\geq r_\alpha$ up to a fixed constant. Increasing the declared profile and selection constants makes the total failure at most the allocation used in the main proof.
\end{proof}

\section{THE OUTER TOURNAMENT}
\label{app:tournament}

We first prove the simultaneous empirical-Bernstein event.

\begin{lemma}[Pairwise empirical Bernstein]
\label{lem:pairwise-appendix}
Let $f_0,\ldots,f_K$ be fixed before an i.i.d. sample of size $m$. For summable weights $\pi_k$ and
\[
 x_{kl}=\log\frac{32}{\delta\pi_k\pi_l},\qquad
 q_{kl}=x_{kl}/m,
\]
define $U_{kl}$ and $\operatorname{rad}_{kl}$ by \eqref{eq:pair-radius}. With probability at least $1-\delta/8$, simultaneously for all pairs,
\begin{align*}
 |(P_m-P)(\ell_{f_k}-\ell_{f_l})|
 &\leq\operatorname{rad}_{kl},\\
 \operatorname{rad}_{kl}
 &\leq3\sqrt{P(f_k\ne f_l)q_{kl}}+5q_{kl}.
\end{align*}
\end{lemma}

\begin{proof}
Fix a pair and let $Z=\ell_{f_k}-\ell_{f_l}\in\{-1,0,1\}$ and $V=Z^2=\one\{f_k\ne f_l\}$. Put $\mu=EZ$, $p=EV$, and let hats denote empirical means. Bernstein's inequality gives, outside an event of probability $2e^{-x}$,
\[
 |\widehat\mu-\mu|
 \leq\sqrt{2px/m}+2x/(3m).
\]
Applying the one-sided inequality to $V$ and using $\sqrt{2pq}\leq p/2+q$ yields
\[
 p\leq2\widehat p+3q=U.
\]
The opposite one-sided inequality gives $U\leq3p+6q$. Substituting the first bound into Bernstein proves coverage by $\sqrt{2Uq}+2q/3$. Substituting the second and using $\sqrt{a+b}\leq\sqrt a+\sqrt b$ proves the displayed upper bound with room in the constants. The pair fails with probability at most $4e^{-x_{kl}}$. Since $\sum_k\pi_k\leq1$, a union over ordered pairs is at most $\delta/8$.
\end{proof}

\begin{proof}[Proof of Lemma~\ref{lem:tournament}]
Apply Lemma~\ref{lem:pairwise-appendix} conditionally on $(S,V)$. If $k$ replaces $c$, then
\[
 R(g_k)-R(g_c)
 \leq\widehat R_W(g_k)-\widehat R_W(g_c)
 +\operatorname{rad}_{kc}<0.
\]
Thus every replacement lowers population risk. If $k$ does not replace $c$, then the empirical difference is at least $-\operatorname{rad}_{kc}$. The same concentration event implies
\[
 R(g_c)-R(g_k)\leq2\operatorname{rad}_{kc}.
\]
The radius upper bound is the second conclusion of Lemma~\ref{lem:pairwise-appendix}.
\end{proof}

For clarity, we give the fixed-point algebra used after challenge $k^*$. Let $a=\Ex(g_c)$, $b=\Ex(g_{k^*})$, and $q=q_{c,k^*}$. Lemma~\ref{lem:regional-appendix} gives
\[
 P(g_c\ne g_{k^*})\leq(C+1)(a^\alpha+b^\alpha).
\]
Failure to replace gives
\begin{equation}
 a-b\leq K\sqrt{(a^\alpha+b^\alpha)q}+Kq.
 \label{eq:appendix-fixedpoint}
\end{equation}
If $a\leq2b$, the oracle bound transfers directly. Otherwise $a-b\geq a/2$ and $b^\alpha\leq a^\alpha$. If the linear $Kq$ term is at least $a/4$, then $a\leq4Kq$. In the other case,
\[
 a\leq K\sqrt{a^\alpha q},
\]
so $a\leq Kq^{1/(2-\alpha)}$. This proves the bound following \eqref{eq:fixedpoint}.

At the first valid constant guess $k^*=\lceil\log_2C\rceil$, the champion index is below $k^*$. Therefore
\[
 x_{c,k^*}\leq K_C+K\log(1/\delta),
\]
with no dependence on $n$. The outer grid eventually contains $k^*$, and every later replacement improves risk. This completes the tournament part of Theorem~\ref{thm:main}.

\section{RECORD PRESERVATION AND SMALL SAMPLES}

We close two endpoint details. For fixed $C_0$ and $\alpha<1$, the profile constants in the preceding construction are bounded over $b\in[0,\alpha]$. The formulas at $b=0$ have finite limits. Let $j_\alpha$ be the last unthinned conservative profile. If its declared radius is a new running minimum, the last retained halving record has radius below $2\rho_{j_\alpha}$. If it is not a new minimum, an earlier retained record is smaller. Hence the oracle record always has radius at most
\[
 2\rho_{j_\alpha}
 \leq4\sup_{0\leq b\leq\alpha}M(C_0,b)r_\alpha.
\]

The proof was written for sufficiently small $s$ and for $N$ large enough that the first valid constant guess appears and conservative profiles do not abort. For fixed $(C,\alpha)$ this excludes only finitely many effective sample sizes $N/D$. Since excess risk is at most one, increasing $K(C,\alpha)$ makes the right side of \eqref{eq:main} at least one throughout this finite range. The theorem therefore holds for every $n$.


\end{document}
